\chapter{From SymPy to Typed IR}\label{ch:translation}

The first stage of the driver checks the symbolic input, sets up the iteration space,
and translates the SymPy assignments into a typed IR block. The classes involved live in
\file{backend/kernelcreation/}. This chapter runs them by hand on a small example.

\section{Kernel analysis}\label{sec:analysis}

\cls{KernelAnalysis} (\file{analysis.py}) walks the assignments once, subexpressions
first and then the main assignments. Within each assignment it visits the right-hand
side before the left-hand side. It enforces the rules of \cref{sec:assignments}:
\begin{itemize}
  \item \textbf{SSA form}: a symbol is assigned at most once, and never after it has been
        read;
  \item \textbf{independence of writes}: a field entry is written at most once, and
        each field (and index) is written at only one offset (switch off with
        \code{allow_double_writes});
  \item \textbf{independence of accesses}: a field read after a write must use the same
        offset as the write (switch off with \code{skip_independence_check});
  \item reduction symbols may only appear on the left of their reduction assignment.
\end{itemize}
Its second job is to collect knowledge. It registers every field in the context with
\code{ctx.add_field(field, const=...)}, and fields that are never written become
\code{const}. Because the right-hand side of an assignment is checked before its own
store is recorded, an in-place update \texttt{f[0,0]}~$\leftarrow$~\texttt{f[1,0]} passes the check
(\cref{sec:constraints}).

\section{Iteration spaces}\label{sec:ispace}

An iteration space (\file{iteration_space.py}) defines the counters of the kernel and
the range of each counter.

\paragraph{Full (dense) iteration spaces.}
A \cls{FullIterationSpace} has one \code{Dimension(start, stop, step, counter)} per
spatial coordinate, with the counters \code{ctr_0}, \code{ctr_1}, \code{ctr_2}. It is
created with
\begin{itemize}
  \item \code{create_with_ghost_layers(ctx, ghost_layers, archetype_field)}, which
        iterates \code{[gl_low, size - gl_high)} in every dimension, or
  \item \code{create_from_slice(ctx, iteration_slice, archetype_field)}, which parses
        the slice with \code{AstFactory.parse_slice} (negative bounds are relative to the
        size).
\end{itemize}
The \term{archetype field} provides the sizes (\code{_size_g_0}, \dots) and the
\code{loop_order}. The driver chooses it with \code{get_archetype_field}: all domain
fields must agree in dimensionality and layout, and the first field in alphabetical
order is used. With \code{ghost_layers=AUTO}, the function \code{infer_ghost_layers}
takes the maximum \code{required_ghost_layers} over all domain field accesses.
\pysnippet{ex17_translation.ispace}
\outsnippet{ex17_translation.ispace}
The loop order \code{(1, 0)} reflects the \code{fzyx} layout: \code{ctr_1} is the outer
loop, and \code{ctr_0} (the contiguous $x$ coordinate) the inner loop.

\paragraph{Sparse iteration spaces.}
A \cls{SparseIterationSpace} iterates over an index list. Its single counter
\code{sparse_idx} runs over the entries of the \code{INDEXED} field, and
\code{get_spatial_counter_declarations} produces the declarations
\code{ctr_i = index_list[sparse_idx].x/y/z} at the top of the loop body (see the sparse
kernel in \cref{sec:sparse}).

\section{Freezing: SymPy to IR}\label{sec:freeze}

\cls{FreezeExpressions} (\file{freeze.py}) is a visitor over SymPy objects. For every
node it looks up a method \code{map_<ClassName>} along the class's method resolution
order and raises a \code{FreezeError} if none exists. The most important rules:

\begin{center}\small
\begin{tabularx}{\linewidth}{@{}l>{\raggedright\arraybackslash}X@{}}
  \toprule
  \textbf{SymPy input} & \textbf{IR output}\\
  \midrule
  \code{Assignment(symbol, rhs)} & \cls{PsDeclaration}\\
  \code{Assignment(field_access, rhs)} & \cls{PsAssignment} to a \cls{PsBufferAcc}\\
  \code{Field.Access} & \cls{PsBufferAcc} with indices \emph{counter + offset} for every spatial dimension, plus the index-dimension entries; \code{PsLookup} for struct members\\
  \code{Symbol}, \code{TypedSymbol} & \cls{PsSymbolExpr} of the context's symbol of that name\\
  \code{Integer}, \code{Float} & untyped \cls{PsConstantExpr}\\
  \code{Rational(p, q)} & \code{p / q} (\cls{PsDiv})\\
  \code{Add} & sum; terms of the form $-1\cdot t$ become \cls{PsSub}\\
  \code{Pow(b, n/d)}, $d\le2$, $|n|\le8$ & \code{sqrt} if $d=2$, then pairwise products; \code{1 / (...)} for $n<0$\\
  other \code{Pow} & \code{pow} (\cls{PsMathFunction})\\
  \code{exp}, \code{sin}, \code{Abs}, \code{Min}, \dots & \cls{PsCall} of a \cls{PsMathFunction}\\
  \code{pi}, \code{E}, \code{oo} & \cls{PsConstantFunction} calls\\
  \code{Piecewise}, \code{ConditionalFieldAccess} & nested \cls{PsTernary}\\
  \code{TypeCast} (\code{tcast}) & \cls{PsCast} (or a typed constant)\\
  relations, \code{And}, \code{Or}, \code{Not} & \cls{PsRel} subclasses, \cls{PsAnd}, \cls{PsOr}, \cls{PsNot}\\
  \code{int_div}, \code{int_rem}, bit operations & \cls{PsIntDiv}, \cls{PsRem}, shifts and bitwise nodes\\
  \code{*ReductionAssignment} & update of a local accumulator \code{<sym>_local}; registered in the context\\
  RNG invocation & \cls{PsRngEngineFunction} call\\
  \bottomrule
\end{tabularx}
\end{center}

Here it is for an expression with powers, a rational, a square root, an exponential and
a free symbol:
\pysnippet{ex17_translation.freeze}
\outsnippet{ex17_translation.sympy}
\irsnippet[nolol=false,caption={After freezing: all constants are still untyped.}]{ex17_translation.frozen}
$f_E^3$ became two multiplications, $1/3$ a division, and $-x/f_C$ a subtraction of
\code{x * (1 / f_C)}. Every field access is now a buffer access with three indices: two
spatial ones and the trivial index dimension of a scalar field.

\section{Typification}\label{sec:typify}

The \cls{Typifier} (\file{typification.py}) assigns a type to every symbol, constant
and expression, and checks that all types fit together. It works with
\term{type contexts} (\cls{TypeContext}). Each expression tree is typified within a
context that has a \term{target type}:
\begin{itemize}
  \item The target type comes from the left-hand side of an assignment, from a typed
        symbol or field access, from a function signature, or from the default for the
        position (index expressions use \code{index_dtype}; conditions use \code{bool}).
  \item Untyped constants simply take the target type (\code{interpret_as}). This is why
        the literal \code{2} becomes \code{2.0f} in a \code{float} kernel.
  \item An untyped symbol on the right-hand side gets \code{default_dtype}. An untyped
        symbol on the left-hand side first tries to infer its type from the right-hand
        side, and otherwise gets \code{default_dtype}.
  \item A typed node whose type differs from the target type is an error. There are no
        implicit conversions, so mixing \code{float32} and \code{float64}, or integers and
        floats, raises a \code{TypificationError} that names the expression and both
        types.
  \item Each operator class declares which types it accepts (numeric, integer or
        boolean operations). Relations yield \code{bool}, or a boolean vector for vector
        operands.
\end{itemize}
\pysnippet{ex17_translation.typify}
\irsnippet[nolol=false,caption={After typification: every constant is \texttt{float32} or \texttt{int64}.}]{ex17_translation.typed}
The free symbol \code{x} was given the context's \code{default_dtype}, here
\code{float32}.

\section{The AST factory}

\cls{AstFactory} (\file{ast_factory.py}) bundles these steps for code that constructs IR:
\begin{itemize}
  \item \code{parse_sympy(obj)}: freeze, then typify, for one expression or assignment;
  \item \code{parse_index(idx)}: parse an expression in an \code{index_dtype} context;
  \item \code{parse_slice(slice, normalize_to)}: turn a Python slice into start, stop
        and step expressions;
  \item \code{loop(ctr_name, slice, body)}, \code{loop_nest(...)},
        \code{loops_from_ispace(ispace, body)}: build \cls{PsLoop} nodes;
  \item \code{axes_cube(...)} and \code{cube_from_ispace(ispace, body)}: build the
        iteration-axes representation used by the next stage (\cref{sec:axes}).
\end{itemize}
Transformations use the factory when they create new code, e.g.\ the remainder loop of a
vectorized kernel. This guarantees that the new nodes are typed consistently.
